\section{张量积}\label{sec:galois-1-8}
读者也许已经接触过实向量空间的``复化''这一概念.下面我们将介绍一种更一般的方法，它使我们能够从环 $A$ 上的模过渡到某个 $A$-代数 $B$ 上的模，其基础是\newterm{galois-tensor-product}{张量积}{tensor product}的概念.

\subsection{双线性映射}\label{sec:galois-1-8-1}
设 $A$ 是交换含幺环，$M,N,P$ 是三个 $A$-模.如果对于每个 $m\in M$ 以及每个 $n\in N$，映射
\begin{align*}
\theta(m,-):N\longrightarrow P,\qquad \theta(-,n):M\longrightarrow P
\end{align*}
都是 $A$-线性的（即 $A$-模同态），则称映射 $M\times N\xrightarrow{\theta}P$ 是 $A$-\newterm{galois-bilinear}{双线性的}{bilinear}.
\begin{example*}
若 $B$ 是 $A$-代数，则乘法映射 $B\times B\longrightarrow B$ 是 $A$-双线性的.
\end{example*}
\begin{example*}
若 $M,N$ 是两个 $A$-模，则映射 $M\times\operatorname{Hom}_A(M,N)\longrightarrow N$，$(m,\varphi)\mapsto\varphi(m)$ 是 $A$-双线性的.
\end{example*}
我们将从 $M\times N$ 到 $P$ 的所有 $A$-双线性映射组成的集合记作 $\operatorname{Bil}_A(M\times N,P)$.它自然地构成一个 $A$-模，其加法为
\begin{align*}
(\theta_1+\theta_2)(m,n):=\theta_1(m,n)+\theta_2(m,n),
\end{align*}
而 $A$ 的作用为 $(a\theta)(m,n):=a\theta(m,n)$.

\subsection{张量积的存在性与唯一性}\label{sec:galois-1-8-2}
\begin{theorem*}
存在一个 $A$-模 $M\otimes_A N$，配备双线性映射
\begin{align*}
M\times N&\longrightarrow M\otimes_A N,\\
(m,n)&\longmapsto m\otimes n,
\end{align*}
并具有如下泛性质：对于任意配备双线性映射 $\theta:M\times N\longrightarrow P$ 的 $A$-模 $P$，存在唯一的 $A$-模同态 $\varphi_\theta:M\otimes N\longrightarrow P$，使得 $\theta(m,n)=\varphi_\theta(m\otimes n)$.换言之，映射 $\varphi\mapsto\varphi\circ(-\otimes-)$ 是一个（$A$-模）同构
\begin{align*}
\operatorname{Hom}_A(M\otimes_A N,P)\xrightarrow{\sim}\operatorname{Bil}_A(M\times N,P).
\end{align*}
此外，配备双线性映射 $(m,n)\mapsto m\otimes n$ 的 $A$-模 $M\otimes_A N$ 在唯一同构的意义下唯一.
\end{theorem*}
\begin{remark*}
这一次，我们通过泛性质而非具体构造来定义一个对象.这是因为，一般而言，这个构造本身并不能告诉我们太多东西.
\end{remark*}
\begin{proof}
与往常一样，唯一同构意义下的唯一性由泛性质推出.为了证明存在性，我们必须``亲手''构造 $M\otimes_A N$，并尽量满足泛性质提出的各项要求.

从自由 $A$-模 $A^{(M\times N)}$ 出发，将其``典范''基记作 $(e_{m,n})_{(m,n)\in M\times N}$.考虑其中由下列元素生成的 $A$-子模 $K$：
\begin{enumerate}
\item $e_{m+m',n}-e_{m,n}-e_{m',n}$ 与 $e_{m,n+n'}-e_{m,n}-e_{m,n'}$，其中 $m,m'\in M$，$n,n'\in N$；
\item $e_{am,n}-ae_{m,n}$ 与 $e_{m,an}-ae_{m,n}$，其中 $m\in M$，$n\in N$，$a\in A$.
\end{enumerate}
于是令 $M\otimes_A N:=A^{(M,N)}/K$，并以 $m\otimes n=\overline{e_{m,n}}$ 表示 $e_{m,n}$ 在此商中的像.\footnote{译注：原文此处及下文写作 $A^{(M,N)}$；按前文的定义，应理解为 $A^{(M\times N)}$.}注意，由 $K$ 的定义，有关系
\begin{align*}
(m+m')\otimes n=\overline{e_{m+m',n}}=\overline{e_{m,n}}+\overline{e_{m',n}}=m\otimes n+m'\otimes n.
\end{align*}
类似地，有
\footnote{译注：此式依原文保留.第二个等号右侧原文仍为 $a\overline{e_{am,n}}$，疑应为 $a\overline{e_{m,n}}$.}
\begin{align*}
(am)\otimes n=\overline{e_{am,n}}=a\overline{e_{am,n}}=a(m\otimes n).
\end{align*}
因此，对每个 $n$，映射 $m\mapsto m\otimes n$ 都是 $A$-线性的.类似地，对每个 $m$，映射 $n\mapsto m\otimes n$ 都是 $A$-线性的.最终，映射 $(m,n)\mapsto m\otimes n$ 是 $A$-双线性的.于是，对任意 $A$-模 $P$，通过复合得到映射
\begin{align*}
\operatorname{Hom}_A(M\otimes_A N,P)&\longrightarrow\operatorname{Bil}_A(M\times N,P),\\
\varphi&\longmapsto\bigl(\theta_\varphi:(m,n)\mapsto\varphi(m\otimes n)\bigr).
\end{align*}
我们需要证明此映射是双射.

为此，构造一个反向映射.从双线性映射 $\theta:M\times N\longrightarrow P$ 出发.根据自由模的泛性质，存在唯一的 $A$-模同态 $\widetilde\varphi_\theta:A^{(M\times N)}\longrightarrow P$，使得 $\widetilde\varphi_\theta(e_{m,n})=\theta(m,n)$.注意到
\begin{align*}
&\widetilde\varphi_\theta(e_{m+m',n}-e_{m,n}-e_{m',n})\\
&\qquad=\theta(m+m',n)-\theta(m,n)-\theta(m',n)=0,
\end{align*}
其中最后一个等号来自 $\theta$ 的双线性.类似地，
\begin{align*}
\widetilde\varphi_\theta(e_{am,n}-ae_{m,n})=\theta(am,n)-a\theta(m,n)=0,
\end{align*}
并且
\begin{align*}
\widetilde\varphi_\theta(e_{m,n+n'}-e_{m,n}-e_{m,n'})=0=\widetilde\varphi(e_{m,an}-ae_{m,n}).
\end{align*}
因此 $K\subset\operatorname{Ker}(\widetilde\varphi_\theta)$.根据商模的泛性质，$\widetilde\varphi_\theta$ 通过唯一的同态
\begin{align*}
\varphi_\theta:A^{(M,N)}/K=M\otimes_A N\longrightarrow P
\end{align*}
分解，此同态将 $m\otimes n=\overline{e_{m,n}}$ 映为 $\widetilde\varphi_\theta(e_{m,n})=\theta(m,n)$.这样就构造了映射
\begin{align*}
\operatorname{Bil}_A(M\times N,P)&\longrightarrow\operatorname{Hom}_A(M\otimes_A N,P),\\
\theta&\longmapsto\bigl(\varphi_\theta:m\otimes n\mapsto\theta(m,n)\bigr).
\end{align*}
显然，这样构造出的两个映射互为逆映射.
\end{proof}
\begin{remark*}[{注意}]
$M\otimes_A N$ 中的一般元素并非形如 $m\otimes n$，而是这类元素的 $A$-线性组合.
\end{remark*}
$M\otimes_A N$ 的元素有时称为\newterm{galois-tensor}{张量}{tensor}.形如 $m\otimes n$ 的元素称为\newterm{galois-elementary-tensor}{初等张量}{elementary tensor}.因此，每个张量都是初等张量的线性组合.

\subsection{函子性质}\label{sec:galois-1-8-3}
设 $M,N$ 如前，$\varphi:M\longrightarrow M'$ 与 $\psi:N\longrightarrow N'$ 是两个 $A$-模同态.
\begin{lemma*}
存在唯一的 $A$-模同态
\begin{align*}
\varphi\otimes\psi:M\otimes_A N\longrightarrow M'\otimes_A N',
\end{align*}
将 $m\otimes n$ 映为 $\varphi(m)\otimes\psi(n)$，其中 $m\in M$，$n\in N$.
\end{lemma*}
\begin{proof}
考虑映射 $\theta:M\times N\longrightarrow M'\otimes_A N'$，它将 $(m,n)$ 映为张量 $\varphi(m)\otimes\psi(n)$.有
\begin{align*}
\theta(am+m',n)&=\varphi(am+m')\otimes\psi(n)\\
&=(a\varphi(m)+\varphi(m'))\otimes\psi(n)\\
&=a(\varphi(m)\otimes\psi(n))+\varphi(m')\otimes\psi(n)\\
&=a\theta(m,n)+\theta(m',n).
\end{align*}
类似地，有 $\theta(m,an+n')=a\theta(m,n)+\theta(m,n')$.因此 $\theta$ 是双线性的，从而存在唯一的同态 $M\otimes_A N\longrightarrow M'\otimes_A N'$，将 $m\otimes n$ 映为 $\theta(m,n)=\varphi(m)\otimes\psi(n)$.
\end{proof}
\begin{remark*}
引理中的唯一性意味着，若 $\varphi':M'\longrightarrow M''$ 与 $\psi':N'\longrightarrow N''$ 是另外两个同态，则有
\footnote{译注：此式依原文保留；其中 $\varphi\circ\psi$ 与 $\psi'\otimes\psi$ 疑分别应为 $\varphi\otimes\psi$ 与 $\psi'\circ\psi$.}
\begin{align*}
(\varphi'\otimes\psi')\circ(\varphi\circ\psi)=(\varphi'\circ\varphi)\otimes(\psi'\otimes\psi).
\end{align*}
\end{remark*}


\subsection{幺半性质}\label{sec:galois-1-8-4}
张量积在 $A$-模的``范畴''上确实扮演着``乘法''的角色.下面的命题说明，它``本质上''是交换的、结合的，对直和满足分配律，并且 $A$ 是一个中性``对象''.
\begin{proposition*}
设 $M,M',N$ 是 $A$-模.
\begin{enumerate}[label=\roman*)]
\item 映射 $M\longrightarrow A\otimes_A M$，$m\mapsto1\otimes m$ 是 $A$-模同构，其逆映射将 $a\otimes m$ 映为 $am$.
\item 存在唯一的 $A$-模同构
\begin{align*}
M\otimes_A N\xrightarrow{\sim}N\otimes_A M,
\end{align*}
将 $m\otimes n$ 映为 $n\otimes m$.
\item 存在唯一的 $A$-模同构
\begin{align*}
(M\oplus M')\otimes_A N\xrightarrow{\sim}(M\otimes_A N)\oplus(M'\otimes_A N),
\end{align*}
将 $(m,m')\otimes n$ 映为 $(m\otimes n,m'\otimes n)$.
\item 存在唯一的 $A$-模同构
\begin{align*}
(M\otimes_A M')\otimes_A N\xrightarrow{\sim}M\otimes_A(M'\otimes_A N),
\end{align*}
将 $(m\otimes m')\otimes n$ 映为 $m\otimes(m'\otimes n)$.
\end{enumerate}
\end{proposition*}
\begin{proof}
\begin{enumerate}[label=\roman*)]
\item 映射 $A\times M\longrightarrow M$，$(a,m)\mapsto am$ 是 $A$-双线性的，因此来自一个同态 $A\otimes_A M\xrightarrow{\varphi}M$，将 $a\otimes m$ 映为 $am$.将命题中的同态记作 $\psi$.有 $\varphi\circ\psi(m)=\varphi(1\otimes m)=m$，故 $\varphi\circ\psi=\operatorname{id}_M$.另一方面，
\begin{align*}
\psi\circ\varphi(a\otimes m)&=\psi(am)=1\otimes(am)=a(1\otimes m)\\
&=(a\cdot1)\otimes m=a\otimes m.
\end{align*}
由于初等张量生成 $A\otimes_A M$，可知 $\psi\circ\varphi=\operatorname{id}_{A\otimes_A M}$.
\item 映射 $M\times N\longrightarrow N\otimes_A M$，$(m,n)\mapsto n\otimes m$ 是 $A$-双线性的，因此来自一个同态 $M\otimes_A N\longrightarrow N\otimes_A M$，将 $m\otimes n$ 映为 $n\otimes m$.由于张量 $m\otimes n$ 生成 $M\otimes_A N$，该性质唯一确定此同态.同样的论证给出一个反向同态，将 $n\otimes m$ 映为 $m\otimes n$.由于张量生成相应的张量积，这样构造的两个同态互为逆同态.
\item 映射 $(M\oplus M')\times N\longrightarrow(M\otimes_A N)\oplus(M'\otimes_A N)$，将 $((m,m'),n)$ 映为 $(m\otimes n,m'\otimes n)$，是 $A$-双线性的，因此来自命题中所述的同态.反过来，映射 $M\times N\longrightarrow(M\oplus M')\otimes_A N$，将 $(m,n)$ 映为 $(m,0)\otimes n$，是双线性的，由此得到同态 $M\otimes_A N\longrightarrow(M\oplus M')\otimes_A N$，将 $m\otimes n$ 映为 $(m,0)\otimes n$.类似地，有同态 $M'\otimes_A N\longrightarrow(M\oplus M')\otimes_A N$，将 $m\otimes n$ 映为 $(0,m')\otimes n$.\footnote{译注：原文此处写作 $m\otimes n$，但像中使用 $m'$，疑应为 $m'\otimes n$.}于是，直和的泛性质给出同态
\begin{align*}
(M\otimes_A N)\oplus(M'\otimes_A N)\longrightarrow(M\oplus M')\otimes_A N,
\end{align*}
它将 $(m\otimes n,m'\otimes n)$ 映为 $(m,0)\otimes n+(0,m')\otimes n=(m,m')\otimes n$.显然，此同态是命题中同态的逆同态.
\item 例如，可以完全仿照 ii) 论证（留给读者）.也可以更优雅地注意到，映射
\begin{align*}
\theta^1:M\times M'\times N&\longrightarrow T^{(1)}:=(M\otimes_A M')\otimes_A N,\\
(m,m',n)&\longmapsto(m\otimes m')\otimes n
\end{align*}
以及
\begin{align*}
\theta^2:M\times M'\times N&\longrightarrow T^{(2)}:=M\otimes_A(M'\otimes_A N),\\
(m,m',n)&\longmapsto m\otimes(m'\otimes n)
\end{align*}
都是 $A$-三线性的，且各自具有如下泛性质：对于任意 $A$-三线性映射 $\theta:M\times M'\times N\longrightarrow P$，存在唯一的 $A$-模同态 $\varphi_\theta^i:T^{(i)}\longrightarrow P$，使得 $\theta=\theta^i\circ\varphi_\theta^i$.\footnote{译注：复合顺序依原文保留；由各映射的定义域与陪域，应为 $\theta=\varphi_\theta^i\circ\theta^i$.}于是，$\varphi_{\theta^2}^1$ 就是命题中的同态，而 $\varphi_{\theta^1}^2$ 是它的逆同构.
\end{enumerate}
\end{proof}
\begin{remark*}
有时，我们将略微滥用记号，把 $M\otimes_A(M'\otimes_A N)$ 与 $(M\otimes_A M')\otimes_A N$ 等同起来，简记为
\begin{align*}
M\otimes_A M'\otimes_A N.
\end{align*}
如上述证明所解释的，它配备了 $A$-三线性映射 $(m,m',n)\mapsto m\otimes m'\otimes n$.
\end{remark*}
\begin{exercise*}
\begin{enumerate}[label=\roman*)]
\item 利用命题的 i) 与 iii)，证明 $A^n\otimes_A M\simeq M^n$.
\item 将命题的 iii) 推广到任意直和，即
\begin{align*}
\left(\bigoplus_{i\in I}M_i\right)\otimes_A N\simeq\bigoplus_{i\in I}(M_i\otimes_A N).
\end{align*}
\item 由此推出，对任意集合 $I$，有 $A^{(I)}\otimes_A M\simeq M^{(I)}$.
\end{enumerate}
\end{exercise*}
下面是一个重要例子，其中张量积具有更明确的形式.

\subsection{自由模的张量积}\label{sec:galois-1-8-5}
\begin{proposition*}
两个自由模的张量积仍是自由模.更准确地说，对于两个集合 $I,J$，将 $e_{i,j}$ 映为 $e_i\otimes e_j$ 的唯一 $A$-模同态
\begin{align*}
\varphi:A^{(I\times J)}\longrightarrow A^{(I)}\otimes A^{(J)}
\end{align*}
是一个同构 $A^{(I\times J)}\xrightarrow{\sim}A^{(I)}\otimes A^{(J)}$.
\end{proposition*}
\begin{proof}
这可以从前一命题结合上述习题推出，不过下面给出详细论证.考虑由
\begin{align*}
\left(\sum_i a_ie_i,\sum_j b_je_j\right)\longmapsto\sum_{i,j}a_ib_je_{i,j}
\end{align*}
定义的映射 $A^{(I)}\times A^{(J)}\longrightarrow A^{(I\times J)}$.可以立即验证它是 $A$-双线性的.因此，存在唯一的 $A$-模同态
\begin{align*}
\psi:A^{(I)}\otimes_A A^{(J)}\longrightarrow A^{(I\times J)},
\end{align*}
将元素 $(\sum_i a_ie_i)\otimes(\sum_j b_je_j)$ 映为 $\sum_{i,j}a_ib_je_{i,j}$，因而特别将 $e_i\otimes e_j$ 映为 $e_{i,j}$.由于 $\psi\circ\varphi(e_{i,j})=e_{i,j}$，有 $\psi\circ\varphi=\operatorname{id}_{A^{(I\times J)}}$.另一方面，对 $m=\sum_i a_ie_i$ 与 $n=\sum_j b_je_j$，有
\begin{align*}
\varphi\circ\psi(m\otimes n)&=\varphi\left(\sum_{i,j}a_ib_je_{i,j}\right)=\sum_{i,j}a_ib_j(e_i\otimes e_j)\\
&=\sum_i a_i\left(e_i\otimes\left(\sum_j b_je_j\right)\right)\\
&=\sum_i a_i(e_i\otimes n)=\left(\sum_i a_ie_i\right)\otimes n=m\otimes n.
\end{align*}
由于初等张量生成 $M\otimes N$（参见定理~\ref{sec:galois-1-8-2} 的构造），这说明 $\varphi\circ\psi=\operatorname{id}_{A^{(I)}\otimes A^{(J)}}$.
\end{proof}
\begin{exercise*}
设 $k$ 是域，$V,W$ 是两个有限维 $k$-向量空间.将 $V$ 的对偶记作 $V^*:=\operatorname{Hom}_k(V,k)$.
\begin{enumerate}[label=\roman*)]
\item 证明存在唯一的同构 $V^*\otimes_k W\xrightarrow{\sim}\operatorname{Hom}_k(V,W)$，将 $f\otimes w$ 映为 $k$-线性映射 $v\mapsto f(v)w$.
\item 设 $W=V$.证明存在唯一的 $k$-线性形式 $V^*\otimes_k V\xrightarrow{\varepsilon}k$，将 $f\otimes v$ 映为 $f(v)$.
\item 证明复合映射
\begin{align*}
\operatorname{End}_k(V)=\operatorname{Hom}_k(V,V)\xrightarrow{\sim}V^*\otimes_k V\xrightarrow{\varepsilon}k
\end{align*}
\begingroup\emergencystretch=3em
就是同态 $\varphi\mapsto\operatorname{tr}(\varphi)$.[可以选取 $V$ 的一组基 $e_1,\ldots,e_n$，将其在 $V^*$ 中的对偶基记作 $e_1^*,\ldots,e_j^*$，\footnote{译注：对偶基末尾的下标 $j$ 依原文保留，疑应为 $n$.}注意到 $e_j^*\otimes e_i$ 对应于 $V$ 的一个自同态，其在所选基下的矩阵是初等矩阵 $E_{ij}$，最后注意到 $\varepsilon(e_j^*\otimes e_i)=\delta_{ij}$.]
\par\endgroup
\end{enumerate}
\end{exercise*}
\begin{remark*}
必须明确指出张量积是在哪个环上取的.例如，考虑两个维数分别为 $d,d'$ 的 $\mathbb C$-向量空间 $V,V'$.那么 $V\otimes_{\mathbb C}V'$ 是维数为 $dd'$ 的 $\mathbb C$-向量空间，因而也是维数为 $2dd'$ 的 $\mathbb R$-向量空间.相比之下，$V\otimes_{\mathbb R}V'$ 是维数为 $(2d)(2d')=4dd'$ 的 $\mathbb R$-向量空间.
\end{remark*}


\subsection{标量扩张}\label{sec:galois-1-8-6}
现在假设 $A$-模 $M$ 是一个 $A$-代数, 并将它记作 $B$. 我们要把 $B\otimes_A N$ 上的 $A$-模结构扩张为 $B$-模结构.
\begin{proposition*}
$B\otimes_A N$ 上存在唯一的 $B$-模结构使得
\begin{align*}
\forall b,b'\in B,\quad\forall n\in N,\qquad b'\cdot(b\otimes n)=(b'b)\otimes n.
\end{align*}
\end{proposition*}
\begin{proof}
我们规定了 $b'$ 在初等张量上的作用, 而这些初等张量生成 $B\otimes_A N$, 所以唯一性随之得到. 还须说明这样规定的 $b'$ 的作用良定义, 并且给出一个 $B$-模结构.

存在性. 记 $\mu_{b'}:B\longrightarrow B$, $b\mapsto b'b$ 为 $B$ 中乘以 $b'$ 的映射. 它是 $B$ 的 $A$-线性自同态. 由\hyperref[sec:galois-1-8-3]{引理 1.8.3}, 存在唯一态射 $\Psi_{b'}:=\mu_{b'}\otimes\mathrm{id}_N:B\otimes_A N\longrightarrow B\otimes_A N$, 将 $b\otimes n$ 映为 $\mu_{b'}(b)\otimes n=(b'b)\otimes n$.

$B$-模结构. 现在要验证映射 $b'\in B\mapsto\Psi_{b'}\in\operatorname{End}_A(B\otimes_A N)$ 是 $A$-代数态射. 等式 $((b'+b'')b)\otimes n=(b'b)\otimes n+(b''b)\otimes n$ 表明 $\Psi_{b'+b''}=\Psi_{b'}+\Psi_{b''}$. 此外,
\begin{align*}
\Psi_{b''b'}(b\otimes n)&=(b''b'b)\otimes n
=\Psi_{b''}((b'b)\otimes n)
=\Psi_{b''}\circ\Psi_{b'}(b\otimes n),
\end{align*}
故 $\Psi_{b''b'}=\Psi_{b''}\circ\Psi_{b'}$.
\end{proof}
上面构造的 $B$-模结构具有如下泛性质.
\begin{proposition*}
对任意 $B$-模 $M$ 和任意 $A$-模态射 $\psi:N\longrightarrow M$, 存在唯一的 $B$-模态射 $\varphi_\psi:B\otimes_A N\longrightarrow M$, 使得 $\varphi(b\otimes n)=b\psi(n)$. 换言之, 映射 $\varphi\mapsto(\psi_\varphi:n\mapsto\varphi(1\otimes n))$ 给出双射
\begin{align*}
\operatorname{Hom}_B(B\otimes_A N,M)\xrightarrow{\sim}\operatorname{Hom}_A(N,M),
\end{align*}
其逆映射为 $\psi\mapsto\varphi_\psi$.
\end{proposition*}
\begin{proof}
仍然由于初等张量 $b\otimes m$ 生成 $B\otimes_A M$, 得到 $\varphi=\varphi_\psi$ 的唯一性. 要得到陈述中的 $A$-模态射 $\varphi$, 只须验证映射 $B\times N\longrightarrow M$, $(b,n)\mapsto bn$ 是 $A$-线性的, 这是显然的. 再验证 $\varphi$ 确为 $B$-模态射, 这由计算
\begin{align*}
\varphi(b'\cdot(b\otimes n))=\varphi((b'b)\otimes n)=b'bn=b'\varphi(b\otimes n)
\end{align*}
得到. 最后, 显然 $\psi_{\varphi_\psi}=\psi$; 另一方面,
\begin{align*}
\varphi_{\psi_\varphi}(b\otimes n)&=b\psi_\varphi(n)=b\varphi(1\otimes n)
=\varphi(b(1\otimes n))=\varphi(b\otimes n).
\end{align*}
这就得到陈述的最后一个断言.\footnote{译注: 本段沿用原文记号. 原文将唯一性论证中的 $N$ 写成 $M$, 并将构造映射写成 $(b,n)\mapsto bn$, 称其为 $A$-线性映射; 按命题中的对象, 相应映射应为 $(b,n)\mapsto b\psi(n)$, 性质应为 $A$-双线性. 计算中的 $bn$ 同样应理解为 $b\psi(n)$.}
\end{proof}
\begin{corollary*}
若 $N$ 是具有基 $(e_i)_{i\in I}$ 的自由 $A$-模, 则 $B\otimes_A N$ 是具有基 $(1\otimes e_i)_{i\in I}$ 的自由 $B$-模.
\end{corollary*}
\begin{proof}
对任意 $A$-模 $M$, 考虑映射 $\varphi\mapsto(\varphi(1\otimes e_i))_{i\in I}$, $\operatorname{Hom}_B(B\otimes_A N,M)\longrightarrow M^I$. 它分解为
\begin{align*}
\operatorname{Hom}_B(B\otimes_A N,M)
&\xrightarrow{\varphi\mapsto\psi_\varphi}\operatorname{Hom}_A(N,M)
\xrightarrow{\psi\mapsto(\psi(e_i))_{i\in I}}M^I.
\end{align*}
命题告诉我们左边的箭头是双射, 而 $(e_i)_{i\in I}$ 是 $N$ 的 $A$-基这一事实告诉我们右边的箭头也是双射. 所以 $(\varphi(1\otimes e_i))_{i\in I}$ 是 $B\otimes_A N$ 的 $B$-基.\footnote{译注: 原文此处写的是任意 $A$-模 $M$, 但 $\operatorname{Hom}_B$ 要求 $M$ 是 $B$-模; 结论中的基也原样写成了 $(\varphi(1\otimes e_i))$, 按陈述应为 $(1\otimes e_i)$.}
\end{proof}
\begin{remark*}
常常将 $1\otimes e_i$ 简记为 $e_i$. 这是一种记号滥用.
\end{remark*}
\begin{example*}
设 $V$ 是一个 $\mathbb R$-向量空间. 在没有张量积时, 常用下列两种方式之一引入 $V$ 的复化 $V_{\mathbb C}$:
\begin{enumerate}
\item 选取 $V$ 的一组 $\mathbb R$-基 $(e_i)_{i\in I}$, 令 $V_{\mathbb C}:=\bigoplus_{i\in I}\mathbb C e_i$;
\item 令 $V_{\mathbb C}:=V\oplus V$, 并定义 $a+ib\in\mathbb C$ 对 $(v,w)$ 的乘法为 $(a+ib)\cdot(v,w):=(av-bw,aw+bv)$.
\end{enumerate}
第一种方法依赖基的选择, 因而不是典范的. 第二种似乎更自然, 但若要推广它, 例如定义一个 $\mathbb Q$-向量空间的``实化'', 就会很困难! 有了张量积, 定义 $V_{\mathbb C}$ 的恰当方式是令 $V_{\mathbb C}:=\mathbb C\otimes_{\mathbb R}V$. 上述推论明确说明了它与第一种方法的联系. 而分解
\begin{align*}
\mathbb C\otimes_{\mathbb R}V
=(\mathbb R\otimes_{\mathbb R}V)\oplus(i\mathbb R\otimes_{\mathbb R}V)
\simeq V\oplus iV
\end{align*}
说明了它与第二种方法的联系.
\end{example*}
\begin{remark*}
设 $M$ 是 $B$-模, $N$ 是 $A$-模. 与\hyperref[sec:galois-1-8-6]{1.8.6} 第一个命题相同的论证表明, $M\otimes_A N$ 上存在唯一的 $B$-模结构, 使 $b\in B$ 对初等张量 $m\otimes n$ 的作用为 $b\cdot(m\otimes n):=(bm)\otimes n$. 换言之, 若 $\rho_b:M\longrightarrow M$ 表示 $M$ 中乘以 $b$ 的映射, 则 $M\otimes_A N$ 中乘以 $b$ 的映射就是 $\rho_b\otimes\mathrm{id}_N$.
\end{remark*}

\subsection{标量扩张的传递性}\label{sec:galois-1-8-7}
先取上面注中的 $M,N$, 再设 $M'$ 是另一个 $B$-模. 下列引理的证明与\hyperref[sec:galois-1-8-4]{命题 1.8.4} 的 iv) 相同.
\begin{lemma*}
存在唯一的 $B$-模同构
\begin{align*}
M\otimes_B(M'\otimes_A N)\xrightarrow{\sim}(M\otimes_B M')\otimes_A N,
\end{align*}
对所有 $m\in M$, $m'\in M'$, $n\in N$, 将 $m\otimes(m'\otimes n)$ 映为 $(m\otimes m')\otimes n$.
\end{lemma*}
取引理中的特例 $M'=B$, 得到同构
\begin{align*}
M\otimes_B(B\otimes_A N)\xrightarrow{\sim}(M\otimes_B B)\otimes_A N\xrightarrow{\sim}M\otimes_A N,
\end{align*}
其中第二个同构为 $i\otimes\mathrm{id}_N$, 而 $i:M\otimes_B B\xrightarrow{\sim}M$ 是\hyperref[sec:galois-1-8-4]{命题 1.8.4} 的 i) 中将 $m\otimes b$ 映为 $bm$ 的同构. 逆同构
\begin{align*}
M\otimes_A N\xrightarrow{\sim}M\otimes_B(B\otimes_A N)
\end{align*}
将 $m\otimes n$ 映为 $m\otimes(1\otimes n)$.

现在假设 $M$ 是一个 $B$-代数, 将它记作 $C$. 因而有两个环态射 $A\longrightarrow B\longrightarrow C$.
\begin{proposition*}
存在唯一的 $C$-模同构
\begin{align*}
C\otimes_A N\xrightarrow{\sim}C\otimes_B(B\otimes_A N),
\end{align*}
将 $c\otimes n$ 映为 $c\otimes(1\otimes n)$; 其逆映射将 $c\otimes(b\otimes n)$ 映为 $cb\otimes n$.
\end{proposition*}
\begin{proof}
我们刚才已经证明, 这些态射存在并且是 $B$-模同构. 只须再说明它们是 $C$-线性的, 而由 $C$ 的作用的定义, 这是显然的.
\end{proof}

\subsection{沿商态射的标量扩张}\label{sec:galois-1-8-8}
这里考虑 $B=A/I$ 的情形, 其中 $I$ 是 $A$ 的理想. 若 $M$ 是 $A$-模, 记
\begin{align*}
IM:=\{m\in M:\ &\exists i_1,\ldots,i_r\in I,\ \exists m_1,\ldots,m_r\in M,\\
&m=i_1m_1+\cdots+i_rm_r\}.
\end{align*}
不难验证这是 $M$ 的一个 $A$-子模 (还可注意到, 若 $M$ 是 $A$ 的理想 $J$, 就得到积理想 $IJ$ 的定义). 按照构造, $I$ 在商 $A$-模 $M/IM$ 上的作用为零. 因而 $A$-模结构 $A\longrightarrow\operatorname{End}_{\mathbb Z}(M/IM)$ 通过 $A/I$ 分解, 从而使 $M/IM$ 成为 $A/I$-模.
\begin{proposition*}
映射 $m\mapsto1\otimes m$ 诱导 $A/I$-模同构
\begin{align*}
M/IM\xrightarrow{\sim}A/I\otimes_A M.
\end{align*}
\end{proposition*}
\begin{proof}
陈述中的映射 $\psi$ 是 $A$-线性的. 它将元素 $im$ 映为 $A/I\otimes_A M$ 中的 $1\otimes(im)=i(1\otimes m)=(i\cdot1)\otimes m$. 但在 $A/I$ 中 $i\cdot1=0$, 故 $\psi(im)=0$. 由于 $IM$ 由形如 $im$ 的元素生成, 得到 $\psi(IM)=0$, 因而该映射通过商 $M/IM$ 分解为 $A$-线性映射 $\bar\psi:M/IM\longrightarrow A/I\otimes_A M$. 反过来, 从投影 $M\xrightarrow{m\mapsto\bar m}M/IM$ 出发. 它属于 $\operatorname{Hom}_A(M,M/IM)$, 标量扩张的泛性质保证存在 $\varphi\in\operatorname{Hom}_{A/I}(A/I\otimes_A M,M/IM)$, 使 $\varphi(\bar a\otimes m)=\bar a\cdot\bar m$. 对所有 $a\in A$, $m\in M$, 有
\begin{align*}
\bar\psi\circ\varphi(\bar a\otimes m)&=1\otimes am=a(1\otimes m)=\bar a\otimes m,\\
\varphi\circ\bar\psi(\bar m)&=\varphi(1\otimes m)=\bar m.
\end{align*}
所以 $\varphi$ 与 $\bar\psi$ 是互逆的双射. 由于 $\varphi$ 是 $A/I$-线性的, 它的逆 $\bar\psi$ 也是如此.
\end{proof}
\begin{example*}
设 $A=\mathbb Z$, $M$ 是一个 Abel $p$-群, $l$ 是不同于 $p$ 的素数. 则 $\mathbb Z/l\mathbb Z\otimes_{\mathbb Z}M=0$. 事实上, 上述推论告诉我们该模同构于 $M/lM$, 而乘以 $l$ 在 $M$ 上是满射: 若 $p^km=0$ 且 $ul+vp^k=1$, 则 $lum=m$.
\end{example*}
\begin{exercise*}
设 $I,J$ 是满足 $I+J=A$ 的两个理想, $M$ 是被 $I\cap J$ 零化的 $A$-模, 即 $(I\cap J)M=0$. 证明 $M=IM\oplus JM$, $IM\simeq M/JM$, $JM\simeq M/IM$.
\end{exercise*}


\subsection{沿局部化的标量扩张}\label{sec:galois-1-8-9}
这里考虑 $B=S^{-1}A$ 的情形, 其中 $S\subset A$ 是乘法子集. 若 $M$ 是 $A$-模, 可以用构造 $S^{-1}A$ 的同样方法构造一个 $S^{-1}A$-模 $S^{-1}M$. 在 $M\times S$ 上赋予等价关系
\begin{align*}
(m,s)\sim(m',s')\quad\Longleftrightarrow\quad\exists t\in S,\ tsm'=ts'm,
\end{align*}
将 $(m,s)$ 在商集 $S^{-1}M$ 中的等价类记为 $\frac ms$. 可以验证, $S^{-1}M$ 上存在唯一的 Abel 群运算使
\begin{align*}
\frac ms+\frac{m'}{s'}=\frac{s'm+s'm}{ss'},
\end{align*}
随后存在唯一的 $S^{-1}A$ 作用使 $\frac at\cdot\frac ms=\frac{am}{ts}$. 称 $S^{-1}M$ 为 $M$ 沿 $S$ 的局部化.\footnote{译注：原PDF此处加法分子印作 $s^{\prime}m+s^{\prime}m$；通常应为 $s^{\prime}m+sm^{\prime}$. 正文保留原式.}
\begin{proposition*}
映射 $M\longrightarrow S^{-1}M$, $m\mapsto\frac m1$ 诱导 $S^{-1}A$-模同构
\begin{align*}
S^{-1}A\otimes_A M\xrightarrow{\sim}S^{-1}M.
\end{align*}
\end{proposition*}
\begin{proof}
陈述中的映射是 $A$-线性的, 因为它将 $am$ 映为 $\frac{am}1=\frac a1\cdot\frac m1$. 由标量扩张的泛性质, 得到一个 $S^{-1}A$-模态射 $\psi:S^{-1}A\otimes_A M\longrightarrow S^{-1}M$. 反过来, 我们希望定义映射 $\varphi$, 将 $\frac ms$ 映为 $\frac1s\otimes m$. 为说明此定义有意义, 设 $(m',s')\sim(m,s)$, 并取 $t$ 使 $tsm'=ts'm$. 有
\begin{align*}
\frac1{s'}\otimes m'&=\frac{ts}{tss'}\otimes m'=ts\left(\frac1{tss'}\otimes m'\right)
=\frac1{tss'}\otimes tsm'\\
&=\frac1{tss'}\otimes ts'm=ts'\left(\frac1{tss'}\otimes m\right)
=\frac{ts'}{tss'}\otimes m=\frac1s\otimes m.
\end{align*}
所以 $\varphi$ 良定义. 验证 $\psi$ 与 $\varphi$ 是互逆的双射留给读者.
\end{proof}
\begin{remark*}
沿用上述记号:
\begin{enumerate}[label=\roman*)]
\item 若 $M$ 的每个元素都被 $S$ 的某个元素零化, 则 $S^{-1}M=0$. 事实上, 若 $tm=0$, 则对任意 $s$ 都有 $\frac ms=\frac{mt}{st}=0$. 例如, 若 $M$ 是有限 Abel 群, 则 $\mathbb Q\otimes_{\mathbb Z}M=0$.
\item 若 $S$ 的每个元素在 $M$ 上的作用都是双射, 则 $S^{-1}M=M$. 事实上, 典范映射 $m\mapsto\frac m1$ 的核是集合 $\{m\in M:\exists t\in S,tm=0\}=\{0\}$. 若将乘以 $s$ 的双射的逆记作 $s_M^{-1}$ (它未必来自 $A$ 中某个元素的作用), 则
\begin{align*}
\frac ms=\frac{ss_M^{-1}(m)}s=\frac{s_M^{-1}(m)}1,
\end{align*}
所以该映射是满射. 例如 $\mathbb Q\otimes_{\mathbb Z}\mathbb Q=\mathbb Q$.
\end{enumerate}
\end{remark*}
\begin{example*}
设 $A=\mathbb Z$, $M$ 是有限 Abel 群. 已知 $M$ 是有限 Abel $p$-子群的积 $\prod_pM_p$, 其中 $p$ 遍历素数 (若尚不知道, 后面会看到). 设 $\mathbb Z_{(p)}$ 是 $\mathbb Z$ 在素理想 $(p)$ 处的局部化, 即沿不同于 $p$ 的素数 $l$ 所生成的乘法子集 $S$ 的局部化. 则
\begin{align*}
\mathbb Z_{(p)}\otimes_{\mathbb Z}M\simeq M_p.
\end{align*}
事实上, 由上注 i), 对 $l\ne p$ 有 $\mathbb Z_{(p)}\otimes M_l=0$, 因为 $M_l$ 被在 $S$ 中可逆的某个幂 $l^k$ 零化. 由上注 ii), 有 $\mathbb Z_{(p)}\otimes M_p=M_p$, 因为所有 $l\ne p$ 都可逆地作用 (逆为乘以 $u$, 其中 $ul+vp^k=1$, 而 $p^k$ 零化 $M_p$).
\end{example*}

\subsection{$A$-代数的张量刻画}\label{sec:galois-1-8-10}
若 $B$ 是 $A$-代数, 我们已注意到乘法 $B\times B\longrightarrow B$ 是 $A$-双线性的. 因而它诱导 $A$-线性映射 $\mu:B\otimes_A B\longrightarrow B$. 反过来, 可以问这样的映射 $\mu$ 应具有哪些性质, 才能成为一个 $A$-代数的乘法.
\begin{proposition*}
设 $\mu:B\otimes_A B\longrightarrow B$ 是 $A$-模态射. 映射 $(b,b')\mapsto b\cdot b':=\mu(b\otimes b')$ 使 $B$ 成为交换 $A$-代数, 当且仅当下列图交换:
\[
\begin{tikzcd}[column sep=large]
B\otimes_A B\otimes_A B\arrow[r,"\mu\otimes\mathrm{id}"]\arrow[d,"\mathrm{id}\otimes\mu"']&B\otimes_A B\arrow[d,"\mu"]\\
B\otimes_A B\arrow[r,"\mu"']&B
\end{tikzcd}
\]
\[
\begin{tikzcd}[column sep=large]
B\otimes_A B\arrow[r,"\mu"]\arrow[d,"{b\otimes b'\mapsto b'\otimes b}"']&B\arrow[d,equal,"\mathrm{id}"]\\
B\otimes_A B\arrow[r,"\mu"]&B
\end{tikzcd}
\qquad
\begin{tikzcd}[column sep=small]
&B\otimes_A B\arrow[dr,"\mu"]&\\
B\arrow[ur,"\mathrm{id}\otimes1"]\arrow[rr,equal,"\mathrm{id}"]&&B
\end{tikzcd}
\]
\end{proposition*}
\begin{proof}
第一个图交换意味着
\begin{align*}
\mu\circ(\mu\otimes\mathrm{id})=\mu\circ(\mathrm{id}\otimes\mu):B\otimes_A B\otimes_A B\longrightarrow B.
\end{align*}
这等价于断言 $\forall b,b',b''\in N$, $(b\cdot b')\cdot b''=b\cdot(b'\cdot b'')$, 即运算 $(b,b')\mapsto\mu(b\otimes b')$ 的结合性.\footnote{译注: 原文此处为 $N$, 按上下文应为 $B$.} 同样, 第二个图交换等价于该运算的交换性, 第三个图交换表示 $1$ 是该运算的单位元. 至于对加法的分配律, 已包含在映射 $(b,b')\mapsto\mu(b\otimes b')$ 的双线性中.
\end{proof}

\subsection{$A$-代数的余积}\label{sec:galois-1-8-11}
这里给定两个 $A$-代数 $B,C$.
\begin{proposition*}
$B\otimes_A C$ 上存在唯一的 $A$-代数结构, 使对所有 $b,b'\in B$, $c,c'\in C$ 都有 $(b\otimes c)(b'\otimes c')=bb'\otimes cc'$.
\end{proposition*}
\begin{proof}
仍然由于初等张量生成 $B\otimes_A C$, 得到乘法的唯一性. 对存在性, 记给出 $B,C$ 中乘法的态射为 $\mu^B:B\otimes_A B\longrightarrow B$, $\mu^C:C\otimes_A C\longrightarrow C$, 并考虑态射
\begin{align*}
\mu^{B\otimes C}:(B\otimes_A C)\otimes(B\otimes_A C)
&\xrightarrow{\sim}(B\otimes_A B)\otimes_A(C\otimes_A C)\\
&\xrightarrow{\mu^B\otimes\mu^C}B\otimes_A C.
\end{align*}
它将 $(b\otimes c)\otimes(b'\otimes c')$ 映为 $(b\otimes b')\otimes(c\otimes c)$, 再映为 $bb'\otimes cc'$.\footnote{译注: 原文中间项确为 $c\otimes c$, 按该映射应为 $c\otimes c'$.} 要说明 $\mu^{B\otimes C}$ 在 $B\otimes_A C$ 上定义一个 $A$-代数结构, 须验证上一个命题中的图交换; 这由 $\mu^B$ 和 $\mu^C$ 的相应图的交换性得到, 验证虽繁琐, 却没有困难.
\end{proof}
\begin{exercise*}
若 $M$ 是 $B$-模, $N$ 是 $C$-模, 证明 $M\otimes_A N$ 上存在唯一的 $B\otimes_A C$-模结构, 使初等张量上的作用为 $(b\otimes c)\cdot(m\otimes n)=bm\otimes cn$.
\end{exercise*}
注意映射 $B\xrightarrow{\mathrm{id}\otimes1}B\otimes_A C$, $b\mapsto b\otimes1$, 以及 $C\xrightarrow{1\otimes\mathrm{id}}B\otimes_A C$, $c\mapsto1\otimes c$, 都是 $A$-代数态射. 配备这两个态射的 $A$-代数 $B\otimes_A C$ 具有下列泛性质.
\begin{proposition*}
设 $D$ 是配备两个代数态射 $\eta:B\longrightarrow D$, $\psi:C\longrightarrow D$ 的 $A$-代数. 则存在唯一代数态射 $B\otimes_A C\xrightarrow{\eta\cdot\psi}D$, 使 $\eta=(\eta\cdot\psi)\circ(\mathrm{id}\otimes1)$, $\psi=(\eta\cdot\psi)\circ(1\otimes\mathrm{id})$. 此外,
\begin{align*}
\eta\cdot\psi=\mu^D\circ(\eta\otimes\psi),
\end{align*}
在初等张量上由 $b\otimes c\mapsto\eta(b)\psi(c)$ 给出. 换言之, 映射 $\theta\mapsto(\theta\circ(\mathrm{id}\otimes1),\theta\circ(1\otimes\mathrm{id}))$ 是双射
\begin{align*}
\operatorname{Hom}_{A\text{-}\mathrm{alg}}(B\otimes_A C,D)\xrightarrow{\sim}
\operatorname{Hom}_{A\text{-}\mathrm{alg}}(B,D)\times\operatorname{Hom}_{A\text{-}\mathrm{alg}}(C,D),
\end{align*}
其逆映射为 $(\eta,\psi)\mapsto\eta\cdot\psi$.
\end{proposition*}
注意它与模的直和的泛性质的相似性; 我们也曾把直和称为``模的余积''.
\begin{proof}
首先直接验证, 由 $\mu^D\circ(\varphi\otimes\psi)$ 定义的 $\varphi\cdot\psi$ 确是代数态射, 并且确有 $\varphi=(\varphi\cdot\psi)\circ(\mathrm{id}\otimes1)$, $\psi=(\varphi\cdot\psi)\circ(1\otimes\mathrm{id})$.\footnote{译注: 证明中的 $\varphi$ 对应命题中的 $\eta$; 原文在此更换了记号.} 为完成证明, 还须说明 $(\theta\circ(\mathrm{id}\otimes1))\cdot(\theta\circ(1\otimes\mathrm{id}))=\theta$. 只须对初等张量 $b\otimes c$ 验证, 而对这样的张量有
\begin{align*}
\theta(b\otimes c)=\theta((b\otimes1)(1\otimes c))=\theta(b\otimes1)\theta(1\otimes c),
\end{align*}
正如所需.
\end{proof}
\begin{example*}
设 $I,J$ 是 $A$ 的两个理想. 有
\begin{align*}
A/I\otimes_A A/J\simeq A/(I+J).
\end{align*}
事实上, 设 $\pi^I:A/I\longrightarrow A/(I+J)$, $\pi^J:A/J\longrightarrow A/(I+J)$ 是典范投影. 由命题, 它们给出态射 $\pi^I\cdot\pi^J:A/I\otimes_A A/J\longrightarrow A/(I+J)$, 将 $(a\bmod I)\otimes(b\bmod J)$ 映为 $ab\bmod(I+J)$. 反过来, 考虑映射 $\varphi:A\longrightarrow A/I\otimes_A A/J$, 将 $a$ 映为 $\varphi(a):=a\cdot((1\bmod I)\otimes(1\bmod J))$. 一方面有 $\varphi(a)=(a\bmod I)\otimes(1\bmod J)$, 所以 $I\subset\operatorname{Ker}(\varphi)$. 另一方面也有 $\varphi(a)=(1\bmod J)\otimes(a\bmod J)$, 所以 $J\subset\operatorname{Ker}(\varphi)$.\footnote{译注: 原文这一式的第一个因子确为 $1\bmod J$, 按张量积所在的模应为 $1\bmod I$.} 因而 $I+J\subset\operatorname{Ker}(\varphi)$, $\varphi$ 通过 $A$-线性态射 $\bar\varphi:A/(I+J)\longrightarrow A/I\otimes_A A/J$ 分解. 按构造,
\begin{align*}
(\pi^I\cdot\pi^J)\circ\bar\varphi(a\bmod(I+J))
&=(\pi^I\cdot\pi^J)(a\cdot((1\bmod I)\otimes(1\bmod J)))\\
&=a\cdot(1\bmod(I+J))=a\bmod(I+J),\\
\bar\varphi\circ(\pi^I\cdot\pi^J)((a\bmod I)\otimes(b\bmod J))
&=\varphi(ab\bmod(I+J))\\
&=ab\cdot(1\bmod I)\otimes(1\bmod J)\\
&=a\cdot(1\bmod I)\otimes(b\bmod J)\\
&=(a\bmod I)\otimes(b\bmod J).
\end{align*}
故 $\bar\varphi$ 与 $\pi^I\cdot\pi^J$ 互逆.\footnote{译注: 原文在第二串等式中将定义域为 $A$ 的 $\varphi$ 用于商类 $ab\bmod(I+J)$; 此处应理解为 $\bar\varphi$.}
\end{example*}
\begin{example*}
将 $X_1$ 映为 $X_1\otimes1$, $X_2$ 映为 $1\otimes X_2$ 的唯一 $A$-代数态射是同构
\begin{align*}
A[X_1,X_2]\xrightarrow{\sim}A[X_1]\otimes_A A[X_2],
\end{align*}
其逆为 $\iota_1\cdot\iota_2$, 其中 $\iota_i$ 是包含映射 $A[X_i]\hookrightarrow A[X_1,X_2]$. 更一般地, 若 $\mathcal N_1,\mathcal N_2$ 是交换幺半群, 则有同构
\begin{align*}
A[\mathcal N_1\times\mathcal N_2]\xrightarrow{\sim}A[\mathcal N_1]\otimes_A A[\mathcal N_2],
\end{align*}
由条件 $e_{(n_1,n_2)}\mapsto e_{n_1}\otimes e_{n_2}$ 确定; 其逆为 $\iota_1\cdot\iota_2$, 其中 $\iota_1:A[\mathcal N_1]\longrightarrow A[\mathcal N_1\times\mathcal N_2]$ 由 $e_{n_1}\mapsto e_{(n_1,0)}$ 确定, 而 $\iota_2:A[\mathcal N_2]\longrightarrow A[\mathcal N_1\times\mathcal N_2]$ 由 $e_{n_2}\mapsto e_{(0,n_2)}$ 确定.
\end{example*}

\begin{example*}[{几何解释}]
上例给出 $\mathbb C$-代数同构 $\mathcal O(\mathbb C^{n_1})\otimes_{\mathbb C}\mathcal O(\mathbb C^{n_2})\xrightarrow{\sim}\mathcal O(\mathbb C^{n_1+n_2})$, 用函数表述, 它将 $f_1\otimes f_2$ 映为函数 $(z_1,z_2)\mapsto f_1(z_1)f_2(z_2)$, 其中 $z_1\in\mathbb C^{n_1}$, $z_2\in\mathbb C^{n_2}$. 更一般地, 若 $V_i\subset\mathbb C^{n_i}$ 是代数子集, 则 $V_1\times V_2$ 是 $\mathbb C^{n_1+n_2}$ 的代数子集, 上述命题给出 $\mathbb C$-代数态射
\begin{align*}
\Psi_{V_1,V_2}:\mathcal O(V_1)\otimes_{\mathbb C}\mathcal O(V_2)\longrightarrow\mathcal O(V_1\times V_2),
\end{align*}
在函数上由同一公式给出. 若将满射 $\mathcal O(\mathbb C^{n_i})\twoheadrightarrow\mathcal O(V_i)$ (由函数限制给出) 记为 $\pi_{V_i}$, 则等式
\begin{align*}
\Psi_{V_1,V_2}\circ(\pi_{V_1}\otimes\pi_{V_2})
=\pi_{V_1\times V_2}\circ\Psi_{\mathbb C^{n_1},\mathbb C^{n_2}}
\end{align*}
表明 $\Psi_{V_1,V_2}$ 是满射. 来证明它也是单射. 为此选取 $\mathcal O(V_2)$ 的一组 $\mathbb C$-基 $(\varepsilon_i)_{i\in I}$. 则元素 $F\in\mathcal O(V_1)\otimes_{\mathbb C}\mathcal O(V_2)$ 唯一地写成 $F=\sum_{i\in I}f_i\otimes\varepsilon_i$, 其中 $f_i\in\mathcal O(V_1)$, 且除有限多个 $i\in I$ 外均为零. 若 $\Psi_{V_1,V_2}(F)=0$, 则对任意 $v_1\in V_1$, 函数 $\sum_{i\in I}f_i(v_1)\varepsilon_i$ 在 $V_2$ 上为零. 因为 $\varepsilon_i$ 构成 $\mathcal O(V_2)$ 的基, 每个 $f_i(v_1)$ 都为零. 由于 $v_1$ 任意, 得到所有 $f_i=0$, 最后得到 $F=0$.
\end{example*}
\begin{remark*}
若从两个代数子集 $V_1,V_2\subset\mathbb C^n$ 出发, 它们的零化理想分别为 $I_1,I_2\subset\mathcal O(\mathbb C^n)$, 则理想 $I_1+I_2$ 零化代数集 $V_1\cap V_2$. 因而上面的第一个例子给出满态射
\begin{align*}
\mathcal O(V_1)\otimes_{\mathcal O(\mathbb C^n)}\mathcal O(V_2)\twoheadrightarrow\mathcal O(V_1\cap V_2).
\end{align*}
可以证明 $V_1\cap V_2$ 的零化理想是 $I_1+I_2$ 的根理想; 这等价于说上述态射的核是 $\mathcal O(V_1)\otimes_{\mathcal O(\mathbb C^n)}\mathcal O(V_2)$ 的幂零根.
\end{remark*}
回到一般记号. 若通过同态 $\mathrm{id}\otimes1$ 将 $B\otimes_A C$ 看作 $B$-代数, 则它具有下列泛性质.
\begin{corollary*}
对任意 $B$-代数 $D$, 映射 $\varphi\mapsto(\psi_\varphi:c\mapsto\varphi(1\otimes c))$ 是双射
\begin{align*}
\operatorname{Hom}_{B\text{-}\mathrm{alg}}(B\otimes_A C,D)\xrightarrow{\sim}\operatorname{Hom}_{A\text{-}\mathrm{alg}}(C,D),
\end{align*}
其逆映射为 $\psi\mapsto\varphi_\psi:b\otimes c\mapsto b\psi(c)$.
\end{corollary*}
\begin{proof}
可用两种方法得到这个结论:
\begin{enumerate}
\item 从上一个命题出发, 将其中的 $\eta$ 固定为给出 $D$ 的 $B$-代数结构的态射.
\item 从\hyperref[sec:galois-1-8-6]{1.8.6} 的第二个命题出发, 取 $M=D$, $N=C$, 并注意双射 $\operatorname{Hom}_B(B\otimes_A C,D)\xrightarrow{\sim}\operatorname{Hom}_A(C,D)$ 将保持乘法的态射映为保持乘法的态射.
\end{enumerate}
\end{proof}
\begin{example*}
将 $X$ 映为 $1\otimes X$ 的唯一 $A$-代数态射是同构
\begin{align*}
A[X]\xrightarrow{\sim}A\otimes_{\mathbb Z}\mathbb Z[X],
\end{align*}
其逆为 $\iota\cdot\kappa$, 其中 $\iota:A\longrightarrow A[X]$ 是典范单射, 而 $\mathbb Z[X]\longrightarrow A[X]$ 是将 $X$ 映为 $X$ 的唯一环态射. 更一般地, 若 $\mathcal N$ 是幺半群, 则有同构
\begin{align*}
A[\mathcal N]\xrightarrow{\sim}A\otimes_{\mathbb Z}\mathbb Z[\mathcal N].
\end{align*}
\end{example*}
\begin{example*}
设 $B$ 是 $A$-代数, $I$ 是 $A$ 的理想, $IB$ 是 $I$ 在 $B$ 中的像所生成的理想. 则有 $B$-代数同构
\begin{align*}
B\otimes_A(A/I)\xrightarrow{\sim}B/IB,
\end{align*}
由 $\pi\cdot\iota$ 给出, 其中 $\pi$ 是投影 $B\longrightarrow B/IB$, 而 $\iota:A/I\longrightarrow B/IB$ 由 $A\longrightarrow B\longrightarrow B/IB$ 通过取商得到. 为构造逆映射, 从 $\mathrm{id}\otimes1:B\longrightarrow B\otimes_A(A/I)$ 出发, 它将 $b$ 映为 $b\otimes(1\bmod I)$. 注意这个 $B$-代数态射将 $ib$ 映为
\begin{align*}
ib\otimes(1\bmod I)=b\otimes(i\bmod I)=b\otimes0=0.
\end{align*}
故它通过态射 $B/IB\longrightarrow B\otimes_A(A/I)$ 分解, 这就是所求的逆.
\end{example*}
\begin{example*}
有同构
\begin{align*}
\varphi:\mathbb C\otimes_{\mathbb R}\mathbb C\xrightarrow{\sim}\mathbb C\times\mathbb C,
\qquad z_1\otimes z_2\mapsto(z_1z_2,z_1\bar z_2).
\end{align*}
为验证它良定义, 可使用命题并注意 $\varphi=\varphi_1\cdot\varphi_2$, 其中 $\varphi_i:\mathbb C\longrightarrow\mathbb C\times\mathbb C$ 由 $\varphi_1(z)=(z,z)$, $\varphi_2(z)=(z,\bar z)$ 定义. 注意 $\varphi(z\otimes1)=(z,z)$, $\varphi(iz\otimes i)=(-z,z)$. 由此推出 $\varphi$ 是满射, 再由维数相等 (在 $\mathbb C$ 或 $\mathbb R$ 上) 得到它是双射. 更确切地说, 它的逆将 $(z,z')$ 映为 $\frac12(z\otimes1-iz\otimes i+z'\otimes1+iz'\otimes i)$.

也可如下理解该同构. 从将 $X$ 映为 $i$ 的 $\mathbb R$-代数同构 $\mathbb R[X]/(X^2+1)\xrightarrow{\sim}\mathbb C$ 出发, 则
\begin{align*}
\mathbb C\otimes_{\mathbb R}\mathbb C
&\simeq\mathbb C\otimes_{\mathbb R}(\mathbb R[X]/(X^2+1))\\
&=\mathbb C\otimes_{\mathbb R}(\mathbb R[X]\otimes_{\mathbb R[X]}(\mathbb R[X]/(X^2+1)))\\
&=(\mathbb C\otimes_{\mathbb R}\mathbb R[X])\otimes_{\mathbb R[X]}(\mathbb R[X]/(X^2+1))\\
&=\mathbb C[X]/(X^2+1)\\
&=\mathbb C[X]/((X-i)(X+i))\\
&=\mathbb C[X]/(X-i)\times\mathbb C[X]/(X+i)\simeq\mathbb C\times\mathbb C.
\end{align*}
第一行使用了前面的例子, 第二行使用了中国剩余定理.\footnote{译注: 原文所说的行次依其排版; 中国剩余定理用于上述最后一步商环的分解.} 若要将这串同构写明, 可以看到它将 $z\otimes1$ 映为 $(z,z)$ (因为每个同构都是 $\mathbb C$-线性的); 将 $z$ 写成 $a+ib$, 则它把 $1\otimes z$ 映为
\begin{align*}
((a+Xb)\bmod(X-i),(a+Xb)\bmod(X+i))=(a+ib,a-ib)=(z,\bar z).
\end{align*}
所以确实重新得到前面的同构.
\end{example*}
\begin{example*}
试推广上例: 从一个在 $\mathbb Q$ 上有限维的域 $K$ 出发, 考察 $\mathbb C$-代数 $\mathbb C\otimes_{\mathbb Q}K$. 将任意域态射 $\sigma:K\longrightarrow\mathbb C$ 称为 $K$ 到 $\mathbb C$ 的嵌入. 注意这种态射自动是 $\mathbb Q$-线性的, 因而就是 $\mathbb Q$-代数态射. 所有这些嵌入的积是一个 $\mathbb Q$-代数态射 $K\longrightarrow\prod_{\sigma:K\hookrightarrow\mathbb C}\mathbb C$. 因为右端是 $\mathbb C$-代数, 上面的推论给出 $\mathbb C$-代数态射
\begin{align*}
(*)\qquad\mathbb C\otimes_{\mathbb Q}K\longrightarrow\prod_{\sigma:K\hookrightarrow\mathbb C}\mathbb C,
\qquad z\otimes x\mapsto(z\sigma(x))_{\sigma:K\hookrightarrow\mathbb C}.
\end{align*}
后面将证明, 域 $K$ 作为 $\mathbb Q$-代数可由一个元素 $\alpha$ 生成 (这样的元素远非唯一). 这意味着将 $X$ 映为 $\alpha$ 的唯一 $\mathbb Q$-代数态射 $\mathbb Q[X]\longrightarrow K$ 是满射. 记其核为 $I$. 因为 $\mathbb Q[X]$ 是主理想环, 存在唯一的首一多项式 $f\in\mathbb Q[X]$, 使 $I=(f)$. 因而有同构
\begin{align*}
\mathbb Q[X]/(f)\xrightarrow{\sim}K,\qquad X\mapsto\alpha.
\end{align*}
由\hyperref[sec:galois-1-4-3]{1.4.3} 的第二个推论, 它表明 $\deg(f)=\dim_{\mathbb Q}(K)=:n$, 且 $(1,\alpha,\ldots,\alpha^{n-1})$ 是 $K$ 的一组 $\mathbb Q$-基. 于是可以计算
\begin{align*}
\mathbb C\otimes_{\mathbb Q}K
&\simeq\mathbb C\otimes_{\mathbb Q}(\mathbb Q[X]/(f))\\
&=\mathbb C\otimes_{\mathbb Q}\mathbb Q[X]\otimes_{\mathbb Q[X]}(\mathbb Q[X]/(f))
=\mathbb C[X]/(f).
\end{align*}
我们还将证明 $f$ 在 $\mathbb C[X]$ 中分解为 $f=(X-\alpha_1)\cdots(X-\alpha_n)$, 其中 $\alpha_1,\ldots,\alpha_n$ 是 $\mathbb C$ 中两两不同的元素. 中国剩余定理保证
\begin{align*}
\mathbb C[X]/(f)=\mathbb C[X]/(X-\alpha_1)\times\cdots\times\mathbb C[X]/(X-\alpha_n)=\mathbb C^n.
\end{align*}
将由此得到的同构
\begin{align*}
(**)\qquad\mathbb C\otimes_{\mathbb Q}K\xrightarrow{\sim}\mathbb C^n
\end{align*}
\begingroup\emergencystretch=3em
写得更明确些. 因为它是 $\mathbb C$-线性的, 它将 $z\otimes1$ 映为 $(z,z,\ldots,z)$. 此外, 我们已知 $(1,\alpha,\ldots,\alpha^{n-1})$ 是 $K$ 的 $\mathbb Q$-基, 故任意 $x\in K$ 唯一地写成 $x=g(\alpha)$, 其中 $g(X)\in\mathbb Q[X]$ 的次数小于 $n$. 可以看到, 上述同构将 $1\otimes x$ 映为 $(g(X)\bmod(X-\alpha_1),\ldots,g(X)\bmod(X-\alpha_n))$, 即 $(g(\alpha_1),g(\alpha_2),\ldots,g(\alpha_n))\in\mathbb C^n$.
\par\endgroup

$(*)$ 与 $(**)$ 有什么关系? 给定 $K=\mathbb Q[X]/(f)$ 到 $\mathbb C$ 的嵌入, 等价于给定 $X=\alpha$ 在 $\mathbb C$ 中的像; 这个像必须是 $f$ 的根, 因而属于 $\{\alpha_1,\ldots,\alpha_n\}$. 所以 $\sigma\mapsto\sigma(\alpha)$ 给出集合 $\{\sigma:K\hookrightarrow\mathbb C\}$ 与 $\{f\text{ 的根}\}$ 之间的双射. 将满足 $\sigma_i(\alpha)=\alpha_i$ 的嵌入记为 $\sigma_i$. 对上述任意 $x=g(\alpha)$, 都有 $\sigma_i(x)=\sigma_i(g(\alpha))=g(\sigma_i(\alpha))=g(\alpha_i)$. 因此将态射 $(*)$ 改写为
\begin{align*}
\mathbb C\otimes_{\mathbb Q}K\longrightarrow\mathbb C^n,\qquad
z\otimes x\mapsto(z\sigma_1(x),\ldots,z\sigma_n(x)),
\end{align*}
就可看到 $(**)$ 与 $(*)$ 是同一态射, 从而推出 $(*)$ 是同构.
\end{example*}
\begin{example*}
再取上例, 但这次计算 $\mathbb R$-代数 $\mathbb R\otimes_{\mathbb Q}K$. 沿用相同记号, 有 $\mathbb R\otimes_{\mathbb Q}K\simeq\mathbb R[X]/(f)$, 但此时 $f$ 不再完全分裂.\footnote{译注：原文在此断言不再分裂，已照录；一般而言，如果所有共轭根均为实数，仍可能在 $\mathbb R$ 上完全分裂.} 设 $f$ 有 $r$ 个实根, $2s$ 个非实根 (每个非实根的共轭也是根, 故其数目为偶数). 有 $n=r+2s$. 将根编号, 使 $\alpha_1,\ldots,\alpha_r$ 是实根, $\alpha_{r+1},\ldots,\alpha_n$ 是非实根, 并且对 $i=1,\ldots,s$ 有 $\alpha_{r+s+i}=\bar\alpha_{r+i}$. 则在 $\mathbb R[X]$ 中有不可约元分解
\footnote{译注: 原书此处的二次式确实没有一次项中的 $X$; 按共轭根相乘, 应为 $X^2-2\operatorname{Re}(\alpha_i)X+|\alpha_i|^2$. 上式保留原文.}
\begin{align*}
f=f_1\cdots f_{r+s},\qquad
f_i=\begin{cases}
X-\alpha_i,&i\le r,\\
X^2-2\operatorname{Re}(\alpha_i)+|\alpha_i|^2,&i>r.
\end{cases}
\end{align*}
由于 $f_i$ 两两互素, 中国剩余定理给出同构
\begin{align*}
\mathbb R\otimes_{\mathbb Q}K\xrightarrow{\sim}\prod_{i=1}^{r+s}\mathbb R[X]/(f_i)=\mathbb R^r\times\mathbb C^s.
\end{align*}
明确地说, 该同构仍由 $y\otimes1\mapsto(y,\ldots,y)$ 以及 $1\otimes x\mapsto(g(\alpha_1),\ldots,g(\alpha_{r+s}))$ 给出, 其中 $x=g(\alpha)$ 且 $\deg(g)<n$. 用嵌入表述, 它将 $1\otimes x$ 映为 $(\sigma_1(x),\ldots,\sigma_{r+s}(x))$ (注意 $i\le r$ 时 $\sigma_i(K)\subset\mathbb R$).

现在注意, 与复共轭复合得到嵌入集合 $\{\sigma:K\hookrightarrow\mathbb C\}$ 的一个对合置换 (即阶为 $2$ 的置换), 而集合 $\Sigma:=\{\sigma_1,\ldots,\sigma_{r+s}\}$ 是嵌入的共轭类的一组代表.\footnote{译注: 原文将对合解释为阶为 $2$; 若所有嵌入均为实嵌入, 该置换实际上是恒等置换, 阶为 $1$.} 有时将上述同构写成
\begin{align*}
\mathbb R\otimes_{\mathbb Q}K\xrightarrow{\sim}\prod_{\sigma\in\Sigma}\mathbb R_\sigma,
\end{align*}
其中约定: 若 $\sigma$ 是实嵌入, 则 $\mathbb R_\sigma=\mathbb R$; 否则 $\mathbb R_\sigma=\mathbb C$.
\end{example*}

